\section*{Введение}
\addcontentsline{toc}{section}{Введение}

В условиях активного использования электронных образовательных ресурсов задача хранения, систематизации и предоставления учебных материалов приобретает особую актуальность. В учебном процессе преподаватели и студенты постоянно работают с лекциями, методическими указаниями, лабораторными заданиями, презентациями и иными файлами, которые должны быть доступны в удобной, структурированной и безопасной форме. При отсутствии единого веб-инструмента материалы нередко хранятся разрозненно: в облачных сервисах, мессенджерах, на личных страницах преподавателей или в локальных каталогах. Это затрудняет поиск нужной информации, усложняет контроль актуальности файлов и не позволяет полноценно разграничивать доступ пользователей.

Целью курсового проекта является разработка веб-приложения для хранения, систематизации и выдачи учебных материалов с учетом роли пользователя. Для достижения данной цели в работе были решены задачи анализа предметной области, выбора стека технологий, проектирования структуры приложения и модели данных, а также реализации механизмов аутентификации, авторизации, проверки прав доступа, CRUD-операций, загрузки и скачивания файлов и формирования статистики.

Объектом исследования является процесс электронного хранения и предоставления учебных материалов в рамках образовательной деятельности. Предметом исследования выступают методы и средства проектирования и разработки веб-приложения, обеспечивающего централизованную работу с такими материалами. Практическая значимость проекта заключается в создании работоспособного программного решения, которое может использоваться как основа для учебного архива материалов кафедры, дисциплины или отдельного курса.